% ===== §19 =====
\section{Participatory Knowing and Epistemic Exclusivity}
\label{p24:sec:participation}

\noindent
The phenomena described in the previous section pose, by the very structure of their givenness, an epistemological question: how is such a culture known, and by whom? The question is not idle, for the answer discloses a mechanism. If the private forms of an intimate culture are given only in the first and second person, then the knowing of them cannot be the knowing appropriate to a third-person object, and the difference is not a matter of degree but of kind. This section draws out that difference, and shows that the peculiar way intimate culture is known---and the peculiar way it resists being known from outside---is not an incidental feature but direct evidence for the mechanism the pillars will name.

\subsection*{Participatory knowing versus objectifying knowledge}

There are two ways a cultural form can be known, and intimate culture admits only the first. The first is \emph{participatory}: one knows the form by standing in the circuit it constitutes, by being one of the two for whom it is a form, by generating and re-enacting it from within. To know an inside joke, in this sense, is to be one of the two people it is funny to; to know a private rite is to be one of the two who perform it and feel its omission as a wrongness. This knowing is not the possession of information about the form; it is participation in the form's ongoing generation, a knowing that is indistinguishable from being a party to the culture known. The second way is \emph{objectifying}: one knows the form as an external object, surveys it, describes it, catalogues its features, explains its function---the knowing an ethnographer brings to a custom, an analyst to a symptom, an outsider to a practice not their own. Macro-cultural forms admit both kinds of knowing; their participants know them from within, and outsiders may know them, with some loss, from without. Intimate culture admits only the first. There is no objectifying knowledge of an intimate culture that is knowledge of \emph{it}, because the object objectifying knowledge would grasp---the charge, the meaning, the ``ours''---exists only in the participatory circuit and is precisely what objectification steps outside of. One can objectify the outward form, the sound of the private word, the motions of the micro-rite; one cannot thereby know the culture, because the culture was never in the outward form but in the circuit the form serves, and the circuit is closed to all but the two.

\subsection*{Exclusivity as evidence, not accident}

The decisive point is that this epistemic exclusivity is not a contingent fact about privacy---not a matter of the couple keeping secrets, or of outsiders lacking access they might in principle be granted---but structural evidence for the mechanism of intimate culture. An outsider cannot come to understand a couple's private forms the way an anthropologist comes to understand a foreign custom, however long the outsider observes, because the barrier is not one of information. The anthropologist's difficulty is that the custom's meaning is hard to gather from outside; enough patient observation reduces the difficulty, because the meaning, though hard to reach, is there in the public practice to be reached. The outsider to an intimate culture faces a difficulty of a different order: the meaning is \emph{not there in the observable form at all}, because it is lodged in the interwoven structure of drive that only the two occupy, and no quantity of observation can supply what is missing, since what is missing is not information about the form but a place in the fantasy that charges it. This is why an inside joke explained to a third party is not thereby made funny to them: the explanation supplies the information and withholds the only thing that mattered, which was never information but position. The exclusivity, then, is evidence. It shows that the meaning of intimate culture is anchored not in a public symbolic system, where more observation would yield more understanding, but in a private coordinate of drive, where observation from outside yields nothing, because there is nothing of the relevant kind outside to observe. The knowability of a culture only from within is the epistemic signature of a culture whose substance lies in the register of a relational Real that the participatory circuit alone touches.

\subsection*{Understanding as recognition, and its failure as objectification}

Read through the Hegelian line, participatory knowing has a further character that objectifying knowledge lacks: it is a knowing that is at the same time a mutual recognition. To know the culture from within is to know oneself in it and to know the other as the one with whom it is shared, so that the understanding of the shared form is inseparable from the recognition of the self and the other whom the form binds. Understanding an intimate culture is thus a moment of the mutual self-knowledge the relation carries on through its exteriorized forms: in reading the shared idiom rightly, each knows again who they are to the other. And the failure that corresponds to this understanding has a name in the same line: it is the \emph{objectification} that treats a living, generated form as a finished thing to be surveyed, the misrecognition that mistakes the residue for the culture and the outward form for the charge. When one party comes to treat the shared culture as an object---to survey it, to hold it at the distance of a thing possessed rather than a process shared---the participatory knowing has lapsed into an objectifying one, and the culture, at that moment, has begun to die in that party's hands, for it lived only in the participation that the objectifying stance has withdrawn. This is the epistemological form of a danger the ethics of the final part will name: that a culture whose being is participatory can be killed by being treated as a thing, and that the treatment as thing is a failure of recognition before it is anything else.

\subsection*{The material mediation of understanding}

The historical-materialist line adds its qualification here, as it did to the account of drive, and for the same reason: to prevent the phenomenology from floating free into a pure account of consciousness. Even the participatory understanding of an intimate culture is materially mediated. The forms in which a couple can generate and re-enact their private culture, and therefore the forms in which they can understand it from within, are conditioned by their material circumstances---by the time they have to share, the spaces they inhabit, the media through which they communicate when apart, the consumer scripts that press upon their sense of what a shared life should look like. The participatory knowing is not a pure meeting of two consciousnesses in a vacuum; it is a knowing carried on in and through a material world that shapes what can be understood and how. This qualification does not diminish the exclusivity or the participatory character just established; it locates them in the material world, and in doing so it prepares the passage, at the end of this part, from the phenomenology of intimate culture to the political economy of its production and capture, where the material mediation of understanding becomes the material colonization of it. For now the epistemology is established: intimate culture is known only from within, its exclusivity is evidence of its anchorage in the relational Real, its understanding is a mutual recognition that objectification kills, and even this knowing is mediated by the material conditions in which the relation is lived.
